\section{\ai{}}
\label{sec:benchmark}

\ai{} is ten frozen research repositories, each of which poses the same problem in a different
domain: improve this repository's \emph{own} training recipe. Table~\ref{tab:tasks} lists them.
They span generative modelling --- reinforcement learning against an aesthetic reward for diffusion
\citep{black2023ddpo} and discrete denoising diffusion for graphs \citep{vignac2023digress} ---
language-model post-training --- direct preference optimization \citep{rafailov2023dpo}, on-policy
distillation from a reasoning teacher \citep{deepseek2025r1}, reinforcement learning for code against
a contamination-controlled benchmark \citep{jain2024livecodebench}, and machine unlearning by negative
preference optimization on a fictitious-forgetting corpus \citep{zhang2024npo,maini2024tofu} ---
reward modelling scored on RewardBench \citep{lambert2024rewardbench}, one-shot pruning at high
sparsity \citep{sun2024wanda,yin2024owl}, weight averaging over a bank of fine-tuned checkpoints
\citep{wortsman2022modelsoups}, and multi-turn agentic RL \citep{wang2025ragen}. Each repository is
taken as its authors released it, with the recipe it commits to left in place as the thing to be
improved. We chose
them for a property the whole design depends on: each ships a concrete training recipe, a frozen
starting artifact, and a metric that can be recomputed reproducibly under a half-day budget on one
GPU. They are ten distinct research problems, not ten settings of one, which is why the ten metrics
cannot be pooled in their raw units and why every aggregate in this paper is a rate over
per-task comparisons rather than an average of scores.

\paragraph{The task contract.} A task hands the agent a repository at a pinned commit, a starting
artifact, and an objective stated as ``improve the repository's recipe'' --- not ``reach score
$x$''. What is held fixed is stated by construction rather than by a checker: the evaluation asset,
the metric, and the frozen evaluator that computes it. What is open is the method. The agent may
rewrite the training code, the data pipeline, the objective, or the optimizer; the one line it may
not cross is the evaluation itself. This is deliberate: fixing the method
would test whether the agent can follow a prescribed recipe, whereas fixing only the outcome and its
measurement tests whether it can find \emph{a} way to improve the system at all --- and, as
\S\ref{sec:analysis} shows, sometimes a way the repository's authors did not intend.

\paragraph{The lifecycle.} Every configuration passes through the same five stages, drawn in
Figure~\ref{fig:pipeline}. \textbf{(1) Exploration.} The agent works in the repository for four
hours on one GPU, free to read, edit, train and query a \emph{fast proxy} metric without limit.
\textbf{(2) Source patch.} Its submission is extracted as a \emph{source-only} diff against the
pinned commit; nothing from the exploration workspace --- no trained weights, no cached state ---
crosses forward. \textbf{(3) Formal replay.} The patch is applied to a fresh container built from
the pinned image digest and retrained under a twelve-hour budget, so the reported artifact is
produced by the recipe, not carried over from exploration. \textbf{(4) Validation.} The run must
yield a usable artifact with up to three retained checkpoints; a configuration that produces none
is recorded as a behaviour failure rather than silently dropped. \textbf{(5) Independent
evaluation.} Each checkpoint is scored by an evaluator that was frozen before any run began, reads
nothing from the agent's workspace, and computes a \emph{final} metric the agent never sees during
exploration.

\paragraph{The visibility boundary.} The separation between the fast proxy the agent may query and
the final metric it may not is the load-bearing property of the design (the vertical line in
Figure~\ref{fig:pipeline}). It is what makes a reported gain reproducible and hard to fake: an agent
can optimize what it can see, but the number that decides its score is computed after the fact, from
a source-only recipe, by a frozen evaluator the agent cannot run. The separation is one of
\emph{access and timing} rather than of sample disjointness --- on some tasks the cheap proxy is
drawn from the same corpus as the final set, since running the full evaluation as a proxy would cost
more than the exploration budget allows --- so what the boundary guarantees is that no agent can
score a candidate under the deciding metric, not that it has never seen a row that metric will use. A high score therefore certifies
that the \emph{recipe} improves the system, not that the agent found the evaluation. The rest of
the paper is about how far that certification actually reaches, and where the metric behind it stops
meaning what it appears to.

\begin{table*}[t]
\centering
\caption{\textbf{The ten \ai{} tasks.} Each freezes a research repository and asks the agent to
improve that repository's own training recipe. \textbf{Start} ($b$) is the frozen checkpoint the
agent is handed and \textbf{Recipe} ($\rshort$) is the repository's committed recipe, both scored
by the same hidden evaluator; arrows give the metric direction. The suite is ten distinct research
problems rather than ten variants of one, which is why the ten metrics are not comparable in their
raw units. Raw metric keys and evaluation-set sizes are in Appendix~\ref{app:extra}; the repositories are cited where they are introduced in \S\ref{sec:benchmark}.}
\label{tab:tasks}
\small
\begin{adjustbox}{max width=\linewidth}
\begin{tabular}{lllcrr}
\toprule
\textbf{Task} & \textbf{Family} & \textbf{Starting artifact} & \textbf{Measure} & \textbf{Start $b$} & \textbf{Recipe $\rshort$} \\
\midrule
DDPO       & Generative       & Stable Diffusion v1.5 + LoRA        & aesthetic score $\uparrow$      & $5.40$   & $5.53$   \\
DiGress    & Generative       & QM9 discrete graph diffusion        & test NLL $\downarrow$           & $78.05$  & $69.57$  \\
DPO        & Preference opt.\ & merged Zephyr/Mistral-7B policy      & IFEval strict $\uparrow$        & $0.397$  & $0.508$  \\
Model Soup & Merging          & 72 CLIP ingredient checkpoints      & ImageNet-V2 top-1 $\uparrow$    & $0.687$  & $0.686$  \\
OPD        & LM post-training & DeepSeek-R1-Distill-Qwen-1.5B       & AIME24/25 @32 $\uparrow$        & $0.252$  & $0.427$  \\
OpenR1     & LM post-training & Qwen2.5-Coder-1.5B-Instruct         & LiveCodeBench pass@1 $\uparrow$ & $0.102$  & $0.133$  \\
NPO        & LM post-training & Llama-3.2-1B-Instruct               & extraction $\downarrow$ / utility $\uparrow$$^{\dagger}$ & \na{}    & $.063/.479$ \\
OWL        & Compression      & OPT-6.7B dense, $70\%$ target sparsity & WikiText-2 ppl $\downarrow$  & $10.86$  & $53.36$  \\
RAGEN      & Agentic RL       & Qwen2.5-3B-Instruct RL policy       & Sokoban solve rate $\uparrow$   & $0.117$  & $0.170$  \\
BTRM       & Reward modeling  & Mistral-7B-Instruct-v0.2 + reward head & RewardBench score $\uparrow$ & \na{}    & $74.57$  \\
\bottomrule
\end{tabular}
\end{adjustbox}
\par\vspace{2pt}
{\footnotesize Model Soup's start $b$ is a strong single-ingredient reference and its shipped
$\rshort$ is the naive uniform average of all 72 ingredients, which sits just below the best
ingredient --- so $\rshort<b$ there is expected, not a regression. $^{\dagger}$NPO reports two components rather than a scalar, and TOFU, the
NPO paper and this task's own contract all forbid collapsing them into one. We therefore score it by
\emph{dominance}: a configuration beats the recipe only when a checkpoint is better on both components
at once. $18$ of $28$ do; $10$ improve one at the other's cost and are counted as not beating it; none
is worse on both (Appendix~\ref{app:unlearning}). It admits no start row because every unlearning run
moves extraction away from the start by construction. BTRM ships no trained scalar reward head to serve
as a start. All are handled explicitly in the results.}
\end{table*}


\subsection{Experimental setup}
\label{sec:setup}

\paragraph{Systems.} We evaluate six systems and treat each \emph{(model, native harness, reasoning
effort)} triple as a distinct unit, because a model cannot be run without a harness and the two
jointly determine behaviour. Three GPT-5.6 variants (Sol, Terra, Luna) run under Codex across six
reasoning-effort levels, \texttt{none} through \texttt{max}; two Claude~5 variants (Opus, Sonnet)
run under Claude Code across five; and Kimi~K3 runs through the Claude Code harness at a single
\texttt{max}-effort setting. Crossed with the ten tasks this gives $180+100+10=290$ configurations.
The grids are complete but not budget-matched across harnesses, so a comparison that spans harnesses
carries the scaffold with it (\S\ref{sec:standings}); the effort sweep of \S\ref{sec:effort} is
therefore confined to the Codex grid, where the harness is held fixed.

\paragraph{References.} Calling a change an improvement requires something to compare against, and
every reference in this paper is measured under the frozen evaluator rather than quoted from a
publication. One of them carries the results. The \emph{shipped recipe} ($\rshort$) is the
repository's committed training configuration run at the step count it commits to, scored by the same
hidden evaluator as every artifact: beating it means beating what the repository actually ships. Every
task has one, so every scored configuration is comparable, and \S\ref{sec:results} is read against it
throughout.

Two further references exist to test that result rather than to produce it. A \emph{budget-matched
replay} ($\rlong$) runs the identical configuration out to the agent's own twelve hours, opening
exactly the two things the patches most often change --- the training horizon and the checkpoint
cadence --- and nothing else, which makes it the closest thing the suite has to an agent-free control.
It is defined wherever extending a training run is meaningful, which is seven of the ten tasks; on the
other three the question does not arise, since weight averaging and one-shot pruning do no training to
extend and agentic RL halts under the recipe's own early-stopping rule. \S\ref{sec:validity} uses it.
The \emph{start} ($b$) is the frozen artifact the agent is handed, scored as delivered;
Table~\ref{tab:tasks} lists it per task as context for how far each recipe already moves, and
\S\ref{sec:results} says why we compute no pooled rate against it.

The two recipe references bracket the agent rather than flatter it. A committed run is short enough to
yield a single checkpoint, so $\rshort$ meets the agent's best of three with one point, which favours
the agent; the twelve-hour replay retains twelve to twenty-nine checkpoints and is reduced to the best
of them, which favours the reference by more than that same margin. \S\ref{sec:ckpt} measures what the
reduction rule is worth.

\paragraph{Metrics and scoring.} Each task declares one final metric, an evaluation asset, and a
direction; the raw metric keys are in Appendix~\ref{app:extra}. Because the ten metrics are
incommensurable --- an aesthetic score, a perplexity, a solve rate --- we never average them, and
report per-task rates and per-task relative changes instead. Every configuration retains up to
three checkpoints from its formal replay and is scored on the best of them under the task's
direction; this best-of-three convention is a property of the harness, not a choice available to
the agent. It is not a neutral bookkeeping detail: it moves the pooled win rate against the shipped
recipe by nearly eight points relative to reporting final checkpoints, so every rate in
\S\ref{sec:results} should be read as a best-of-three rate. \S\ref{sec:ckpt} gives the mechanism. Missing values are
preserved as missing throughout: a checkpoint that produced no score is never counted as a zero.
